\documentstyle[% 


righttag% 


,11pt% 


,amssymb% 


,russian%               remove in Eng. 


,izvru%                 Set izven% in Eng. 


]{amsart} 


 


%


 


\newcommand{\supp}{\operatorname{supp}} 


\newcommand{\tts}[1]{}


\renewcommand{\baselinestretch}{1.05}


 


%\hoffset=-15mm%                  For Eng. 


\setcounter{page}{23} 


\god{1999} 


\nomer{7~(446)} %                        Remove in Eng. 


%\nomer{Vol.\,43, No.\,7}%               Set in Eng. 


%\str{\pageref{begin}--\pageref{end}}%  Set in Eng. 


\udk{517.982} 


 


\author{\it В.Л.\,Крепкогорский} 


% NB:   ^^ remove in Eng. 


\title{Интерполяция и теоремы вложения для квазинормированных  


пространств Бесова} 


 


\date{} 


% \pagestyle{myheadings}%         Set in Eng. 


\pagestyle{plain} %              Remove in Eng. 


\begin{document} 


% \thispagestyle{plain}%          Set in Eng. 


\maketitle 


\label{begin} 


 


Применим функтор $(\cdot,\cdot)_{\theta,q}$ к паре пространств Бесова 


$B^{s_i}_{p_i}$. 


Результат хорошо известен [1] в двух случаях: а)~при выбранном 


$\theta$ параметр $q$ принимает одно определенное значение 


$q=q_\theta:=(1-\theta)/p_0+\theta/p_1$; 


б)~параметр $q\in(0,\infty]$ может быть любым, но $p_0=p_1$. 


Получающиеся в результате интерполяционные пространства были 


найдены Питре в 60-х годах [2]. 


 


При дополнительных условиях $s_i{=}1/p_i+b$, $i=0,1$, 


$\forall q{\in}(0,\infty]$ пространства 


$(B^{s_0}_{p_0}(T),B^{s_1}_{p_1}(T))_{\theta,q}$ ($T$ --- 


единичная окружность на $C$) 


изучались в [3] и [4]. Интерполяция для 


пар $(BMO,B)$ в случае $q=q_\theta:=\theta/q_1$ исследовалась в 


[5]. Наконец, сводку результатов по интерполяции в парах 


общего вида $(B,B)$ и $(BMO,B)$ при $q=q_\theta$ можно найти в 


([6], п.\,3.9.4). 


 


В [7] дано описание пространств 


$(B^{s_0}_{p_0},B^{s_1}_{p_1})_{\theta,q}$ в случае, 


когда $1<p_i<\infty$, $-\infty<s_i<\infty$, $\forall q\in(0,\infty]$. 


Нормы пространств Бесова 


задаются при этом с помощью последовательности сверток. Переход 


к общему случаю $0<p_i<\infty$ затруднен тем обстоятельством, что 


при $p_i<1$ имеет место $l^{s_i}_{p_i}[L_{p_i}]\notin S'$, и обычно 


используемая конструкция 


ретракции не имеет смысла. В данной работе получено 


описание пространств 


$(B^{s_0}_{p_0},B^{s_1}_{p_1})_{\theta,q}=BL^{s,k}_{p,q}$ 


в общем случае. Показано, 


что пространства типа $BL$ наследуют многие свойства пространств 


Бесова. В частности, они образуют систему пространств, 


``замкнутую относительно теорем вложения''. 


 


\section{Основные определения и обозначения} 


 


Квазинормой назовем функционал $\|\cdot\|$, обладающий обычными 


свойствами нормы, за исключением неравенства треугольника, 


которое заменено более слабым неравенством 


$\|f+g\|\le\gamma\{\|f\|+\|g\|\}$, 


где постоянная $\gamma>1$. 


 


Назовем квазинормированной (банаховой) решеткой квазинормированное 


(банахово) пространство функций $E$, квазинорма (норма)  


которого обладает свойством монотонности 


$$ 


|f(x)|\le|g(x)|\ \; \forall x\Rightarrow\|\,f\mid E\,\|\le 


\|\,g\mid E\,\|. 


$$ 


 


Пусть даны две квазинормированные решетки $E(X)$ и $G(Y)$, а 


также функция двух переменных $f(x,y)$, $x\in X$, $y\in Y$. Смешанной 


квазинормой 


назовем функционал 


$\|\,f\mid G[E]\,\|=\big\|\,\|\,f(x,y)\mid E(X)\,\|\,\mid G(Y)\big\|$. 


 


Пусть $F$ --- оператор Фурье, $S$ --- пространство Шварца и $S'$ --- 


пространство умеренных распределений на $R_n$. Обозначим через 


$\Phi(R_n)$ 


совокупность всех систем функций 


$\varphi=\{\varphi_j(x)\}^\infty_{j=0}\subset S(R_n)$ таких, 


что 


\begin{align*} 


\supp\varphi_0&\subset \{x:|x|\le 2\};\\ 


\supp\varphi_j&\subset\{x:2^{j-1}\le|x|\le 2^{j+1}\}, 


\ \; j=1,2,3,\dotsc; 


\end{align*} 


для каждого мультииндекса $\alpha$ существует положительное число 


$c_\alpha$, 


при котором 


$$ 


2^{j|\alpha|}|D^\alpha\varphi_j(x)|\le c_\alpha 


$$ 


для всех $j=0,1,2,\dotsc$ и всех $x\in R_n$ и 


$\sum\limits^\infty_{j=0}\varphi_j(x)=1$ для каждого $x\in R_n$. 


 


Пусть $l^s_q$ --- пространство числовых последовательностей; 


$L_{p,q}(X,\omega,\mu)$ --- пространство Лоренца с весом $\omega$ на 


множестве $X$ с 


мерой $\mu$; 


$B^s_{p,q,(r)}(R_n)$ --- пространство Бесова, $B^s_p:=B^s_{p,p,(p)}$; 


$F^s_{p,q,(r)}(R_n)$ --- пространство Лизоркина--Трибеля с 


квазинормами 


\begin{gather*} 


\|(a_j)^\infty_{j=0}\mid l^s_q\|=\|a_j2^{js}\mid l_q\|;\\ 


\|f\mid L_{p,q}\|=\bigg(\int^\infty_0(t^{1/p}(f)^{*(\mu)} 


(t))^q\frac{dt}{t}\bigg)^{1/q}, 


\end{gather*} 


где $f^{*(\mu)}$ --- невозрастающая равноизмеримая относительно меры 


$\mu$ 


перестановка функции $f$; 


\begin{alignat*}{2} 


\|f\mid B^s_{p,q,(r)}\|&=\|F^{-1}\varphi_j F f\mid 


l^s_q[L_{p,r}]\|, &\quad (\varphi_j)&\in\Phi;\\ 


\|f\mid F^s_{p,q,(r)}\|&=\|F^{-1}\varphi_j F f\mid 


L_{p,r}[l^s_q]\|, &\quad (\varphi_j)&\in\Phi. 


\end{alignat*} 


Здесь предполагается, что $-\infty<s<\infty$, $0<p<\infty$, 


$0<q,r\le\infty$. 


 


При интерполяции пар пространств Бесова получаются пространства 


$$ 


BL^{s,k}_{p,q}:=\{f\in S':\|f\mid BL\|= 


\|(F^{-1}\varphi_j F f)\mid L^{s,k}_{p,q}\|<\infty\}, 


$$ 


где $L^{s,k}_{p,q}:=L_{p,q}(2^{j(s-k/p)},\ m_n\times(2^{jk}\nu))$ --- 


пространство Лоренца $L_{p,q}$, 


состоящее из функций, определенных на $R_n\times N$, 


$N=\{0,1,2,3,\dotsc\}$, 


с весом $2^{j(s-k/p)}$, $m_n$ --- лебегова мера на $R_n$, $2^{jk}\nu$ 


--- атомическая 


мера на $N$ с мерой атома $\{j\}$, равной $2^{jk}$, $k$ --- угловой 


коэффициент 


прямой, проходящей через точки $(1/p_0,s_0)$ и $(1/p_1,s_1)$. 


 


\section{Интерполяция в классе квазинормированных пространств  


Бесова} 


 


В этом разделе доказана интерполяционная теорема для 


пространств $B^s_p(R_n)$, аналогичная теореме 1 из [7]. Отличие 


состоит 


в том, что здесь будет рассмотрен общий случай,  


когда $0<p<\infty$. При этом возникают трудности, связанные с тем, 


что  


обычно оператор ретракции не определен на пространстве $l^s_p[L_p]$,  


$p<1$. Здесь приведено доказательство, основанное на том, что 


для $K$-метода достаточно доказать действие ретракции не на всем  


пространстве $l^s_p[L_p]$, а на более узком множестве.


 


\begin{thm} 


Пусть $-\infty<s_0,s_1<\infty$, $0<\theta<1$, $0<p_0,p_1<\infty$, 


$p_0\ne p_1$, $0<q\le\infty$, 


$k$ --- угловой коэффициент прямой, проходящей через точки 


$(1/p_0,s_0)$ 


и $(1/p_1,s_1)$. Если $s=(1-\theta)s_0+\theta s_1$ и 


$1/p=(1-\theta)/p_0+\theta/p_1$, то 


$$ 


(B^{s_0}_{p_0}(R_n),B^{s_1}_{p_1}(R_n))_{\theta,q}=BL^{s,k}_{p,q}. 


$$ 


\end{thm} 


 


{\bf Доказательство} разобьем на несколько частей. 


 


1) Проинтерполируем связанные с пространствами Бесова 


квазинормированные 


структуры. Рассматривая структуры $l^s_p[L_p]$ как пространства $L_p$ 


на $R^n\times N$ с весом $2^{js}$, $j\in N$, применим формулу Фрейтага 


[8]. 


Подробные вычисления можно найти в [7]. В результате получаем равенства 


\begin{gather*} 


(l^{s_0}_{p_0}[L_{p_0}],l^{s_1}_{p_1}[L_{p_1}])_{\theta,q}= 


L^{s,k}_{p,q}=L_{p,q}(2^{j(s-k/p)},\ m_n\times(2^{jk}\nu)), \\ 


%


\|(g_j)\mid L^{s,k}_{p,q}\|=\bigg(\int^\infty_0 


(t^{1/p}(g_j(x)2^{j(s-k/p)})^{*(\mu)})^q\frac{dt}{t}\bigg)^{1/q},\quad 


\mu=m_n\times(2^{jk}\nu). 


\end{gather*} 


 


2) Пусть $H$ --- подмножество квазинормированной структуры $A$. 


Через $M(H)$ обозначим множество 


$\{f\in A;\exists g\in H;\ |g|\ge|f|\}$. 


 


Пусть $L_0$ и $L_1$ --- квазинормированные структуры, $B_0$ и $B_1$ --- 


квазинормированные пространства, $T\in\cal L(B_i,L_i)$, причем 


$\|f\mid B_i\|=\|T f\mid L_i\|$, $i=0,1$. 


Через $T(H)$ обозначим образ множества $H$ при отображении 


$T$. Допустим, что существует линейный оператор $R$, определенный на  


$M(T(B_0+B_1))$, ограниченный на 


$L_i\cap M(T(B_0+B_1))$, $i=0,1$, и $RT=I$. 


Сравним $K$-функционалы $K(t,f,B_0,B_1)$ и $K(t,Tf,L_0,L_1)$. 


Очевидно, 


\begin{multline*} 


K(t,f,B_0,B_1)=\inf_{f_0+f_1=f}(\|Tf_0\mid L_0\|+ 


t\|Tf_1\mid L_1\|)\ge \\


\ge \inf_{g_0+g_1=Tf} 


(\|g_0\mid L_0\|+t\|g_1\mid L_1\|)=K(t,Tf,L_0,L_1). 


\end{multline*}  


С другой стороны, инфимум в последнем $K$-функционале достигается 


на элементах $g_0,g_1\in M(T(B_0+B_1))$, поэтому оператор $R$ определен 


на них и $Rg_0+Rg_1=R(Tf)=f$. Очевидно, 


$\|Rg_0\mid B_0\|+t\|Rg_1\mid B_1\|\le C(\|g_0\mid L_0\|+ 


t\|g_1\mid L_1\|)$. 


Следовательно, $K(t,f,B_0,B_1)\le 


C\inf\limits_{g_0+g_1=Tf}(\|Rg_0\mid B_0\|+t\|Rg_1\mid B_1\|\le 


C K(t,Tf,L_0,L_1)$. 


Итак, в этом случае 


\begin{equation} 


K(t,f,B_0,B_1)\simeq K(t,Tf,L_0,L_1)\quad 


\text{и} 


\quad \|f\mid(B_0,B_1)_{\theta,q}\|\simeq 


\|Tf\mid(L_0,L_1)_{\theta,q}\|. 


\end{equation} 


 


3) Рассмотрим любую пару пространств 


$(B^{s_0}_{p_0},B^{s_1}_{p_1})$. 


Определим 


операторы $Tf:=\{F^{-1}\varphi_j Ff\}^\infty_{j=0}$; 


$R(\{f_j\}^\infty_{j=0})=\sum\limits^\infty_{S'j=0} 


F^{-1}\varphi_j Ff_j$, $\varphi_j\in\Phi$. 


Покажем, что оператор $R$ определен на 


$M(T(B^{s_0}_{p_0},B^{s_1}_{p_1}))$. 


Для этого 


используем теорему вложения 2.7.1 [1], в которой утверждается, 


что $B^s_{p,q}(R_n)\subset B^{s'}_{p',q}(R_n)$ при $0<p_0\le 


p_1\le\infty$, $0<q\le\infty$, $-\infty<s'\le s<\infty$ 


и $s-n/p=s'-n/p'$. Очевидно также, $B^s_{p,q}\subset B^{s^*}_{p,q}$ 


при $s^*<s$, 


$B^s_{p,q}\subset B^s_{p,q^*}$ при $q<q^*$. Поэтому 


$B^s_p=B^s_{p,p}\subset B^{s'}_{p',p}\subset B^{s'}_{p',p'}= 


B^{s'}_{p'}$ 


при $s'\le s-n/p+n/p'$, $p\le p'$. Для произвольной пары пространств 


$(B^{s_0}_{p_0},B^{s_1}_{p_1})$ 


положим $s=s_0$, $p=p_0$ или $s=s_1$, $p=p_1$, а $p'=\max(p_i,2)$, 


$s'=\min(s_i-n/p_i+n/p')$, $i=0,1$. Тогда 


$B^{s_i}_{p_i}\subset B^{s'}_{p'}$ и 


$1<p'<\infty$. Более того, 


$B^{s_0}_{p_0}+B^{s_1}_{p_1}\subset B^{s'}_{p'}$, и 


$T(B^{s_0}_{p_0}+B^{s_1}_{p_1})\subset l^{s'}_{p'}[L_{p'}]$, 


$M(T(B^{s_0}_{p_0}+B^{s_1}_{p_1}))\subset l^{s'}_{p'}[L_{p'}]$. 


Поэтому оператор $R$ определен на 


$M(T(B^{s_0}_{p_0}+B^{s_1}_{p_1}))$. 


Как видно из п.\,2), 


теперь достаточно убедиться, что оператор $R$ ограничен 


на $l^{s_i}_{p_i}[L_{p_i}]\cap M(T 


(B^{s_0}_{p_0}+B^{s_1}_{p_1}))$. 


 


4) Докажем ограниченность оператора $R$, считая $s=s_0$, $p=p_0$ или  


$s=s_1$, $p=p_1$. Для доказательства воспользуемся теоремой 1.6.3 [1]. 


\medskip 


 


{\bf Теорема А.}{\it Пусть


$0<p<\infty$, $0<q\le\infty$, $\Omega=\{\Omega_k\}^\infty_{k=0}$ 


--- последовательность 


компактных подобластей в $R_n$ и $d_k>0$ --- диаметр $\Omega_k$. 


Если $\varkappa>n/2+n/\min(p,q)$, то существует постоянная $C$ 


такая, что неравенство 


$$ 


\|F^{-1}M_k Ff_k\mid L_p(l_q)\|\le C\sup_\ell 


\|M_\ell(d_\ell\cdot)\mid H^\varkappa_2\|\, 


\|f_k\mid L_p(l_q)\| 


$$ 


выполняется для всех систем 


$\{f_k\}^\infty_{k=0}\in L^\Omega_p(l_q)$ и для всех 


последовательностей 


$\{M_k(x)\}^\infty_{k=0}\subset H^\varkappa_2$. 


Здесь 


$$ 


\{f_k\}^\infty_{k=0}\in L^\Omega_p(l_q)\Rightarrow 


f_k\in L^{\Omega_k}_p=\{f\in L_p:\supp Ff\subset \Omega_k;\  


\|f\mid L^{\Omega_k}_p\|=\|f\mid L_p\|<\infty\}. 


$$ 


} 


\medskip 


 


Выберем значения параметров $q$, $\varkappa$, $r$, положив $q:=p$, 


$\varkappa:=n/2+n/p$, 


$r$ --- любое число, большее $\varkappa$. Зададим последовательность 


$\{f_j\}\in M(T(B^{s_0}_{p_0}+B^{s_1}_{p_1}))$ 


и оценим норму 


\begin{multline*} 


\|R(f_j)\mid B^s_p\|=\Big\|\Big(F^{-1}\varphi_k F\Big(\, 


\sum^\infty_{j=0}F^{-1}\varphi_j Ff_j\Big)\Big)^\infty_{k=0} 


\,\Big|\,l^s_p[L_p]\Big\|= \\


%


=\Big\|\Big(F^{-1}2^{ks}\varphi_k F 


\Big(\,\sum^\infty_{j=0}F^{-1}\varphi_j Ff_j\Big) 


\Big)^\infty_{k=0}\,\Big|\,L_p[l_p]\Big\|= \\ 


% 


=\Big\|\Big(\,\sum^\infty_{j=0}F^{-1}2^{ks} 


\varphi_k\varphi_j Ff_j\Big)^\infty_{k=0}\,\Big|\, 


L_p[l_p]\Big\|\le \|(F^{-1}2^{ks}\varphi_k\varphi_{k-1} 


Ff_{k-1})\mid L_p[l_p]\|+ \\ 


% 


+\|(F^{-1}2^{ks}\varphi_k\varphi_k Ff_k)\mid 


L_p[l_p]\|+\|(F^{-1}2^{ks}\varphi_k\varphi_{k+1}Ff_{k+1} 


)\mid L_p[l_p]\|. 


\end{multline*} 


Оценим, используя теорему А, 


$$ 


\|(F^{-1}2^{ks}\varphi_k\varphi_{k+\varepsilon} 


Ff_{k+\varepsilon})^\infty_{k=\max(-\varepsilon,0)}\mid 


L_p[l_p]\|,\quad \varepsilon=-1,0,1. 


$$ 


Для того чтобы $f_{k+\varepsilon}\in L^{\Omega_k}_p[l_p]$, нужно 


положить $\Omega_k=\{x\in R_n:|x|\le 2^{k+\varepsilon+1}\}$ и 


$d_k=2^{k+\varepsilon+1}$. Тогда 


\begin{multline*} 


\|(F^{-1}2^{ks}\varphi_k\varphi_{k+\varepsilon} 


Ff_{k+\varepsilon})\mid L_p[l_p]\|\le 


%%


C\sup_k\|\varphi_k(d_k\cdot)\varphi_{k+\varepsilon} 


(d_k\cdot)\mid H^\varkappa_2\|\,\|(f_{k+\varepsilon})2^{ks}\mid 


L_p[l_p]\|\le \\ 


\le C\sup_k\|\varphi_k(2^{k+\varepsilon+1}x)\varphi_{k+\varepsilon} 


(2^{k+\varepsilon+1}x)\mid W^r_2\|\,\|(f_k)^\infty_{k=0} 


\mid L_p[l^s_p]\|. 


\end{multline*} 


Оценим 


\begin{multline*} 


\|\varphi_k(2^{k+\varepsilon+1}x)\varphi_{k+\varepsilon} 


(2^{k+\varepsilon+1}x)\mid W^r_2\|=\Big(\,\sum_{|\alpha|\le r} 


\|D^\alpha(\varphi_k(2^{k+\varepsilon+1}x) 


\varphi_{k+\varepsilon}(2^{k+\varepsilon+1}x))\mid 


L_2\|\Big)^{1/2} \le \\ 


\le C\Big(\,\sum_{|\beta|+|\gamma|\le r}\|D^\beta 


(\varphi_k(2^{k+\varepsilon+1}x))D^\gamma(\varphi_{k+\varepsilon} 


(2^{k+\varepsilon+1}x))\mid L_2\|^2\Big)^{1/2}. 


\end{multline*} 


Положим $2^{k+\varepsilon+1}x=u$, тогда 


$$ 


|D^\beta_x(\varphi_k(2^{k+\varepsilon+1}x))|= 


|2^{(k+\varepsilon+1)|\beta|}D^\beta_u 


(\varphi_k(u))|\le C.


$$ 


Аналогично показывается, что 


$|D^\gamma_x(\varphi_{k+\varepsilon}(2^{k+\varepsilon+1}x))|\le C^*$. 


При этом 


постоянные $C$ и $C^*$ не зависят от $k$. Кроме того, 


$\varphi_k(2^{k+\varepsilon+1}x)=0$ и 


$\varphi_{k+\varepsilon}(2^{k+\varepsilon+1}x)=0$, если $|x|>2$. 


Следовательно, независимо от $k$ 


$$ 


\|D^\alpha\varphi_k(2^{k+\varepsilon+1}x) 


\varphi_{k+\varepsilon}(2^{k+\varepsilon+1}x)\mid L_2\|\le C. 


$$ 


Поэтому $\sup\limits_k\|\varphi_k(d_k\cdot)\varphi_{k+\varepsilon} 


(d_k\cdot)\mid H^\varkappa_2\|<\infty$, и мы получаем, что оператор 


$R$ ограничен на 


$l^{s_i}_{p_i}[L_{p_i}]\cap M(T 


(B^{s_0}_{p_0}+B^{s_1}_{p_1}))$. 


В силу (1) это означает 


$$ 


\|f\mid(B^{s_0}_{p_0},B^{s_1}_{p_1})_{\theta,q}\|\simeq 


\|T f\mid(l^{s_0}_{p_0}[L_{p_0}],l^{s_1}_{p_1}[L_{p_1}])_{\theta,q}\| 


=\|T f\mid L^{s,k}_{p,q}\|=\|f\mid BL^{s,k}_{p,q}\|.\quad\square 


$$ 


 


{\bf Замечание.} Из равенства (1) и дальнейшего доказательства 


следует


$$ 


K(t,f,B^{s_0}_{p_0},B^{s_1}_{p_1})\simeq K(t,F^{-1} 


\varphi_j Ff,l^{s_0}_{p_0}[L_{p_0}],l^{s_1}_{p_1}[L_{p_1}]). 


$$ 


 


Можно применить полученный результат для описания интерполяции 


в паре пространств $BMO(R_n)$ и $B^s_p(R_n)$. 


Определим пространство 


$BMO(R_n):=\Big\{f\in L^{\text{loc}}_1(R_n):\|f\mid BMO\|= 


\sup\limits_Q\Big(\frac{1}{|Q|}\int\limits_Q|f-f_Q|dx\Big)<\infty 


\Big\}$, 


где $Q$ --- $n$-мерный куб и $f_Q:=\frac{1}{|Q|}\int\limits_Q f\,dx$. 


 


В [5] доказано, что 


$$ 


(BMO,\,B^{s_1}_{p_1,q_1,(r_1)})_{\theta,q}= 


B^s_{p,q,(q)} 


$$ 


при $s_1\ne 0$, $1/p=\theta/p_1$, $1/q=\theta/q_1$, 


$s=\theta s_1$, $p_1,q_1,r_1\in(0,\infty)$. 


В частности, полагая $q_1=p_1=r_1$, получим формулу 


$$ 


(BMO,\,B^{s_1}_{p_1})_{\theta,q}=B^s_p. 


$$ 


Заметим, что эти интерполяционные формулы позволяют найти 


пространство $(\cdot,\cdot)_{\theta,q}$, только если $q=\theta/q_1$, 


т.\,е. только при 


одном значении из всего интервала $(0,\infty)$. Следующая теорема 


позволяет сделать это при любом $q\in(0,\infty]$. 


 


\begin{thm} 


Пусть $s_1\ne0$, $1/p=\theta/p_1$, $s=\theta s_1$, $p_1\in(0,\infty)$, 


$k=p_1s_1$, $q\in(0,\infty]$, $\theta\in(0,1)$. 


Тогда 


$$ 


(BMO,\,B^{s_1}_{p_1})_{\theta,q}=BL^{s,k}_{p,q}. 


$$ 


\end{thm} 


 


\begin{pf} 


Построим систему координат на плоскости, в 


которой пространства $B^s_p$ будем изображать точками $(1/p,s)$. В 


этом 


случае соотношения $1/p=\theta/p_1$, $s=\theta s_1$ означают, что 


точка 


$(1/p,s)$ лежит на отрезке, соединяющем точки $0(0,0)$ и 


$A(1/p_1,s_1)$, 


причем делит его в отношении $(1-\theta){:}\theta$. Можно сказать, что 


с точки 


зрения теории интерполяции пространство $BMO$ соответствует 


пространству $B^0_\infty$. Выберем параметры $\theta_0$ и $\theta_2$ 


так, чтобы $0<\theta_0<\theta<\theta_2<1$. 


По теореме о реитерации при $1/p_0=\theta_0/p_1$, 


$1/p_2=\theta_2/p_1$, $s_0=\theta_0s_1$, $s_2=\theta_2s_1$ 


и соответствующем $\lambda\in(0,1)$ получим равенство 


$$ 


(BMO,\,B^{s_1}_{p_1})_{\theta,q}=((BMO,\, 


B^{s_1}_{p_1})_{\theta_0,p_0},\,(BMO,\, 


B^{s_1}_{p_1})_{\theta_2,p_2})_{\lambda,q}= 


(B^{s_0}_{p_0},B^{s_2}_{p_2})_{\lambda,q}=BL^{s,k}_{p,q}, 


$$ 


где $k$ --- угловой коэффициент прямой, проходящей через точки 


$A_0(1/p_0,s_0)$ и $A_2(1/p_2,s_2)$. На этой же прямой лежат точки 


$0(0,0)$ 


и $A_1(1/p_1,s_1)$. Легко видеть, что ее угловой коэффициент 


$k=s_1p_1$. 


\end{pf} 


 


\section{Теоремы вложения для пространств $BL$} 


 


Пространства типа $BL$ ``наследуют'' ряд свойств, характерных 


для пространств Бесова. Например, для них можно доказать теоремы 


вложения. Интерполируя известные результаты о вложениях 


пространств $B(R_n)$, получим следующую теорему. 


 


\begin{thm} 


Пусть $0<p_0\le p_1<\infty$, $0<q\le\infty$, $-\infty<s_1<s_0<\infty$, 


$-\infty<k_0,k_1<\infty$. Тогда 


\begin{itemize} 


\item[а)] $BL^{s_0,k_0}_{p_0,q_0}(R_n)\subset 


BL^{s_1,k_1}_{p_1,q_1}(R_n)$ 


при $s_0>s_1-n/p_1+n/p_0;$ 


\item[б)] $BL^{s_0,n}_{p_0,q_0}(R_n)\subset 


BL^{s_1,n}_{p_0,q_0}(R_n)$ 


при $s_0-n/p_0=s_1-n/p;$ 


\item[в)] $B^{s_0}_{p_0}(R_n)\subset 


BL^{s_1,k_1}_{p_1,q_1}(R_n);$ 


$BL^{s_0,k_0}_{p_0,q_0}(R_n)\subset 


B^{s_1}_{p_1}(R_n)$ 


при $s_0>s_1-n/p_1+n/p_0;$ 


\item[г)] $BL^{m+n/p,n}_{p,q}\subset C^m(R_n)$, $0<p\le\infty$, 


$0<q<1$, $m=0,1,2,\dotsc$ 


\item[д)] $BL^{s,k}_{p,q}(R_n)\subset C^m(R_n)$ 


при $s>m+n/p;$ 


\item[е)] $BL^{s_0,k}_{p,q}(R_n)\subset\cal R^s(R_n):= 


B^s_{\infty,\infty}(R_n)$, $s>0$, $0<p<\infty$, $0<q\le\infty$, 


$s_0>s+n/p$. 


\end{itemize} 


\end{thm} 


 


Представим пространство $R_n$ как декартово произведение 


$R_n=R_{n-1}\times R_1$, 


$x\in R_n\Rightarrow x=\{x',x_n\}$, 


$x'=\{x_1,x_2,x_3,\dotsc,x_{n-1}\}$. Известно, что 


для функций из пространства Бесова $B^s_p(R_n)$ имеют смысл следы 


$f(x',0)$, $\frac{\partial f}{\partial x_n}(x',0),\dotsc, 


\frac{\partial^r f}{\partial x^r_n}(x',0)$ 


при условии, что 


\begin{equation} 


s>r+1/p\,\max(0,(n-1)(1/p-1)). 


\end{equation} 


 


Для любого пространства $BL^{s,k}_{p,q}(R_n)$ такого, что справедливо 


неравенство (2), можно найти такое пространство $B^{s'}_{p'}$, для 


которого 


выполняется условие (2) и 


$BL^{s,k}_{p,q}(R_n)\subset B^s_p(R_n)$. В этом случае 


функции из $BL^{s,k}_{p,q}(R_n)$ имеют следы на $R_{n-1}$, 


принадлежащие $S'$. 


Определим на $BL^{s,k}_{p,q}(R_n)$ оператор 


$$ 


\Re f:=\Big\{f(x',0),\frac{\partial f}{x_n}(x',0),\dotsc, 


\frac{\partial^r f}{\partial x^r_n}(x',0)\Big\}. 


$$ 


 


Рассмотрим декартово произведение квазинормированных пространств 


$\prod\limits^r_{j=0}B_j$. Определим на нем квазинорму 


$\Big\|(f_j)^r_{j=0}\,\Big|\,\prod\limits^r_{j=0}B_j\Big\|:= 


\sum\limits^r_{j=0}\|f_j\mid B_j\|$. 


 


\begin{thm} 


Если $0<p<\infty$, $0<q<\infty$, $s>r+1/p\,\max(0,(n-1)(1/p-1))$, то 


$\Re$ --- ретракция 


$$ 


BL^{s,k}_{p,q}(R_n)\quad 


\text{на} 


\quad 


\prod^r_{j=0}BL^{s-1/p-j,k-1}_{p,q}(R_{n-1}). 


$$ 


\end{thm} 


 


\begin{pf} 


Так как в теореме 2.7.2 [1] указано, что 


оператор $\Re$ является ретракцией $BL^s_p(R_n)$ на 


$\prod\limits^r_{j=0}B^{s-1/p-j}_{p,q}(R_{n-1})$, то 


очевидно, оператор $\Re$ -- ретракция 


$(BL^s_p(R_n))_{\theta,q}$ на 


$\Big(\prod\limits^r_{j=0}B^{s-1/p-j}_{p,q}(R_{n-1})\Big)_{\theta,q}$. 


На произведении $\prod\limits^r_{j=0}B_j$ зададим нормы 


$$ 


\Big\|(g_j)\,\Big|\,\prod^r_{j=0}B_j\Big\|:=\sum\|(g_j)\mid B_j\|. 


$$ 


Заметим, что $\sum\|(g_j)\mid B_j\|\simeq\max\|(g_j)\mid B_j\|$. Для 


любого набора $(B^0_j,B^1_j)$ очевидно, 


$$ 


K\Big(t,(g_j),\prod^r_{j=0}B^0_j,\ \prod^r_{j=0}B^1_j\Big)= 


\sum^r_{j=0}K(t,g_j,B_j)\simeq \max K(t,g_j,B_j). 


$$ 


Тогда 


\begin{multline*} 


\Phi_{\theta,q}(K(t,(g_j),B^0_j,B^1_j))= 


\Phi_{\theta,q}\Big(\,\sum^\infty_{j=0}K(t,(g_j),B^0_j,B^1_j)\Big) 


\le \\ 


\le \sum^r_{j=0}\Phi_{\theta,q}(K(t,(g_j),B^0_j,B^1_j))\simeq 


\max\Phi_{\theta,q}(K(t,(g_j),B^0_j,B^1_j))\le \\ 


%


\le \Phi_{\theta,q}(\max K(t,g_j,B^0_j,B^1_j))\le \Phi_{\theta,q} 


\Big(\,\sum^r_{j=0}K(t,g_j,B^0_j,B^1_j)\Big). 


\end{multline*} 


Отсюда следует 


$$ 


\Phi_{\theta,q}\Big(K\Big(t,(g_j),\prod^r_{j=0}B^0_j,\  


\prod^r_{j=0}B^1_j\Big)\Big)\simeq\sum^r_{j=0} 


\Phi_{\theta,q}(K(t,g_j,B^0_j,B^1_j))= 


\Big\|(g_j)\,\Big|\,\prod^r_{j=0}(B_j)_{\theta,q}\Big\|. 


$$ 


 


Выберем пару точек $(1/p_i,s_i)$, $i=0,1$, так, чтобы, во-первых, 


$s>r+1/p\,\max(0,(n-1)(1/p-1))$, во-вторых, точка $(1/p,s)$ 


должна лежать 


в середине отрезка с концами в точках $(1/p_i,s_i)$, наконец, 


в-третьих, пусть угловой коэффициент прямой, проходящей через 


точки $(1/p_i,s_i)$, будет равен $k$. В этом случае 


$k=(s_1-s_0)/(1/p_1-1/p_0)$. 


Подсчет углового коэффициента прямой, проходящей через точки 


$(1/p_i,s_i-1/p_i-j)$, приводит к результату $k_j=k-1$. Поэтому 


оператор $\Re$ отображает 


$BL^{s,k}_{p,q}(R_n)= 


(B^{s_0}_{p_0},B^{s_1}_{p_1})_{1/2,q}$ 


на 


$$ 


\prod^r_{j=0}B^{s-1/p-j,k-1}_{p,q}(R_{n-1})= 


\Big(\,\prod^r_{j=0}B^{s_0-1/p_0-j,k-1}_{p_0}(R_{n-1}), 


\ \prod^r_{j=0}B^{s_1-1/p-j,k-1}_{p_1}(R_{n-1})\Big)_{1/2,q} 


$$ 


и является ретракцией. 


\end{pf} 


 


\begin{thebibliography}{99} 


 


\bibitem{8} 


Трибель X. {\em Теория функциональных пространств}. -- М.: Мир, 1986. 


-- 447 с. 


\bibitem{1} 


Peetre J. {\em Sur les spaces de Besov} // Compt. Rend. Acad. Sci. -- 


Paris, 


1967. -- V.\,264. -- \No\,6. -- P.\,281--283. 


\bibitem{2} 


Пеллер В.В. {\em Описание операторов Ганкеля класса $\sigma_p$ при 


$p>0$, исследование скорости рациональной аппроксимации и другие 


приложения} // Матем. сб. -- 1983. -- Т.\,122. -- \No\,4. -- 


С.\,481--510.  


\bibitem{3} 


Пекарский А.А. {\em Классы аналитических функций, определенные 


наилучшими рациональными приближениями в $H_p$} // 


Матем. сб. -- 1985. -- Т.\,127. -- \No\,1. -- C.\,3--20. 


\bibitem{4} 


Peetre J., Svensson E. {\em On the generalized Hardy's inequality 


of McGehee, Pigno and Smith and the problem of interpolation 


between BMO and a Besov space} // Math. Scand. -- 1984. -- V.\,54. --  


P.\,221--241. 


\bibitem{5} 


Bruddnyi Yu.A., Krugliak N.Ya. {\em Interpolation functors and 


interpolation spaces}. I. -- North-Holland, 1991. - 717 р. 


\bibitem{6} 


Крепкогорский В.Л. {\em Интерполяция в пространствах 


Лизоркина--Трибеля 


и Бесова} // Матем. сб. -- 1994. -- Т.\,185. -- \No\,7. -- С.\,63--76. 


\bibitem{7} 


Freitag D. {\em Real interpolation of weighted $L_p$-spaces} // Math. 


Nachr. - 1978. -- Bd.\,86. -- S.\,15--18. 


 


\end{thebibliography} 


 


\vskip 2cm 


 


\post{\it Казанское высшее военное}{\it Поступила} 


\post{\it командно-инженерное училище}{19.03.1996} 


 


\label{end} 


\end{document} 


 


